\section{Уроки 10--11: ArUco-метки и визуальная обратная связь}

\subsection{Цели уроков}
\begin{itemize}
\item Детектировать ArUco-метки в видеопотоке.
\item Оценивать положение маркера через \texttt{solvePnP}.
\item Управлять дроном на основе визуальных признаков.
\end{itemize}

\subsection{Опорные листинги}
\lstinputlisting[style=GeoskanStyle,language=pythoncustom,caption={Базовая детекция ArUco}]{../code/python/examples/aruco_examples/detect_aruco.py}
\lstinputlisting[style=GeoskanStyle,language=pythoncustom,caption={Координаты и оси маркера}]{../code/python/examples/aruco_examples/detect_aruco_coordinates.py}
\lstinputlisting[style=GeoskanStyle,language=pythoncustom,caption={Полёт с визуальной стабилизацией по маркеру}]{../code/python/examples/aruco_examples/aruco_flight.py}

\subsection{Логика регулирования}
\begin{itemize}
\item По дальности до маркера корректируется скорость \texttt{vy}.
\item По положению центра маркера в кадре корректируется \texttt{yaw\_rate}.
\item При потере маркера выполняется остановка команды движения.
\end{itemize}

\begin{figure}[H]
\centering
\begin{tikzpicture}
\draw[GeoskanBlue,thick] (0,0) rectangle (7,4);
\draw[dashed] (2.3,0) -- (2.3,4);
\draw[dashed] (4.7,0) -- (4.7,4);
\node at (1.15,3.5) {Поворот влево};
\node at (3.5,3.5) {Норма};
\node at (5.85,3.5) {Поворот вправо};
\fill[GeoskanOrange] (5.5,2.0) circle (2.5pt);
\end{tikzpicture}
\caption{Разделение кадра для управления рысканьем}
\end{figure}

\begin{taskbox}[Минимальный набор упражнений]
\begin{enumerate}
\item Настройте пороги дальности 1.0--1.5 м под свою площадку.
\item Измените размер метки в \texttt{size\_of\_marker} и оцените влияние на координаты.
\item Добавьте вывод текущего режима: \texttt{approach}, \texttt{hold}, \texttt{retreat}.
\end{enumerate}
\end{taskbox}

\subsection{Критерии успешности}
\begin{enumerate}
\item Маркер обнаруживается устойчиво в 8 из 10 попыток.
\item Команды поворота соответствуют положению маркера в кадре.
\item При потере маркера дрон останавливается без дрейфа.
\end{enumerate}
